\chapter{Conclusion}
\label{zaver}
This thesis introduced a~deterministic and repeatable method for evaluating privacy-{\linebreak}enhancing tools that block network requests. The~proposed method preserves the~structure of network requests, enabling the~assessment of transitive blocking effects. The~method was implemented, validated through unit testing, and then applied to selected content-blocking tools.

The evaluation used data from 937~visited web pages, capturing network requests and DNS responses. Potential fingerprinting was also detected using the~\textit{JShelter Fingerprint Detector}. Request structures were reconstructed as directed trees, linking each resource to any associated fingerprinting activity. A~\textit{BIND 9} DNS server was employed to replay the~original DNS responses using custom-generated zone files. The~tools were tested via a~simulation server replaying recorded requests. Blocked requests were logged, and blocking outcomes were propagated through the~request trees to compute metrics.

The blocking performance of three browsers, \emph{Avast Secure Browser}, \emph{Brave}, and \emph{Firefox}, was evaluated along with five content-blocking extensions: \emph{Adblock Plus}, \emph{Ghostery}, \emph{uBlock Origin}, \emph{uBlock Origin Lite} and \emph{Privacy Badger}. These extensions were tested in their Firefox and Chrome versions when available; in cases where only one version existed, only that version was tested. Ghostery for Firefox delivered the~best overall blocking effectiveness, achieving the~highest scores in most evaluation metrics. uBlock Origin and uBlock Origin Lite also performed well, with results similar to but slightly below those of Ghostery for Firefox. Although the~evaluation included Privacy Badger, its results are unreliable due to the~limitations of the~request simulation process. All tools were tested in their default state, meaning the~results might significantly change if a~user changes their configuration. Notably, many tested tools blocked over \emph{20}~\% of all network requests, underscoring the~significance of using content blockers to reduce unnecessary or potentially invasive web traffic.

Several limitations of the~proposed approach were identified. The~most significant limitation is that the~structure of real-world web requests does not always conform to a~directed tree. In some cases, the~requests form a~general graph with loops, affecting the~evaluation's accuracy as the~proposed approach solves the~loops by introducing artificial duplicates. Another limitation is that the~simulation does not replicate the~original environment, which limits the~effectiveness of evaluating tools that rely on heuristics or dynamic blocking rules. Future work could address these limitations, refine the~evaluation me\-tho\-do\-lo\-gy, expand the~set of metrics, and improve the~evaluation system's overall performance.
